The polarization fraction was selected from the region showing evidence of self-scattering (uniform minor axis alignment). Three methods were tested to determine the most appropriate polarization fraction for this region:
\begin{enumerate}
    \item \POLFsiNS: The polarization fraction at the pixel with maximum \SIns.
    \item $\mathcal{P}_{\mathrm{F, max\,\POLIeq}}$: The polarization fraction at the pixel with maximum polarized intensity.
    \item  $\mathcal{P}_{\mathrm{F, Gaussian}}$: The mean polarization fraction within the uniform region of the Gaussian model, defined as the region where $W_{\text{uniform}}$ is within 1\% of the peak value of 1.
\end{enumerate}
\noindent 

\noindent The values from these three methods at each wavelength are summarized in Table \ref{tab: polf_comparison}. The best representation of the observations was Method 1, where the polarization fraction was taken at the pixel with the maximum \SI value \SC{You need to justify this} . 
% This selected value is referred to as the observed polarization fraction and is denoted by \POLFobsNS. Note that the polarization fractions used here are referred to as observed to distinguish them from the model, but they are debiased. 




% ----------------------------------------------
% ----------------------------------------------
\begin{table}[h!]
\centering
\caption{Polarization fraction values from various methods tested in the region of observed dust self-scattering.}
\label{tab:polf_comparison}
\begin{tabular}{lllll}
\toprule
\textbf{Parameter} & \textbf{\bfour} & \textbf{\bfive} & \textbf{\bsix} & \textbf{\bseven} \\
\midrule

$\mathcal{P}_{\mathrm{F, max\,I}}$
& 1.37
& 1.40
& 1.46
& 1.04 \\

$\mathcal{P}_{\mathrm{F, max\,POLI}}$
& 1.56
& 1.51
& 1.49
& 1.08 \\


$\mathcal{P}_{\mathrm{F, Gaussian}}$
& 1.28
& 1.19
& 1.44
& 0.86 \\

\bottomrule
\end{tabular}
\label{tab: polf_comparison}
\end{table}
% ----------------------------------------------
% ----------------------------------------------